% ============================================================================
%  PROBLEMA DE LA INVESTIGACIÓN
%  Análisis térmico de cables eléctricos subterráneos
%  Luis Koc - Herbert Melendez  ---  MIA 3er Ciclo A 2026-1 - FIIS/UNMSM
%  Mayo 2026
% ============================================================================
\documentclass[aspectratio=169]{beamer}
\usetheme{Madrid}
\usecolortheme{seahorse}
\usefonttheme{professionalfonts}

\usepackage[utf8]{inputenc}
\usepackage[T1]{fontenc}
\usepackage[spanish,es-nodecimaldot]{babel}
\usepackage{lmodern}
% Fix babel-Spanish \es@chf error with fonts that have \hyphenchar=-1
\makeatletter
\def\es@chf{\ifnum\hyphenchar\font>\m@ne\char\hyphenchar\font\fi}
\makeatother
\usepackage{ragged2e}
\usepackage{booktabs}
\usepackage{colortbl}
\usepackage{xcolor}
\usepackage{array}
\usepackage{tabularx}
\usepackage{enumitem}
\usepackage{hyperref}
\usepackage{amsmath}

\usepackage{tikz}
\usetikzlibrary{
  mindmap,
  trees,
  shapes.geometric,
  arrows.meta,
  positioning,
  backgrounds,
  fit,
  calc,
  decorations.pathmorphing
}

% Carpeta con figuras reutilizadas del beamer propuesta
\graphicspath{{c:/usr/ths_mia_fiis/tesis_MIA_lk/temas_cursos_Tesis_1/graficos_ppt/}}

% ---- Beamer settings -------------------------------------------------------
\setbeamertemplate{navigation symbols}{}
\setbeamertemplate{footline}[frame number]
\setbeamertemplate{bibliography item}[text]

% ---- Comandos de referencia ------------------------------------------------
\newcommand{\refcite}[1]{{\footnotesize\color{gray!70!black} #1}}
\newcommand{\rcite}[1]{{\scriptsize\color{gray!70!black} #1}}

% ---- Separadores de sección ------------------------------------------------
\AtBeginSection[]{%
  \begin{frame}[plain, noframenumbering]
    \vfill
    \centering
    \begin{beamercolorbox}[sep=12pt, center, shadow=true, rounded=true]{title}
      \usebeamerfont{title}\insertsectionhead\par%
    \end{beamercolorbox}
    \vfill
  \end{frame}
}

% ---- Portada ---------------------------------------------------------------
\title{El Problema de la Investigación}
\subtitle{Análisis térmico de cables eléctricos subterráneos:\\
          identificación de la brecha y estado del arte}
\author{Luis Koc \and Herbert Meléndez}
\institute{Maestría en Inteligencia Artificial --- FIIS/UNMSM\\
           Trabajo de Investigación I $\cdot$ MIA 3er Ciclo A 2026-1 $\cdot$ Grupo G91}
\date{Mayo 2026}

% ============================================================================
\begin{document}
% ============================================================================

% ---------------------------------------------------------------------------
% PORTADA
% ---------------------------------------------------------------------------
\begin{frame}
  \titlepage
\end{frame}

% ---------------------------------------------------------------------------
% ESTRUCTURA
% ---------------------------------------------------------------------------
\begin{frame}{Estructura de la presentación}
  \begin{columns}[t]
    \column{0.48\textwidth}
    \tableofcontents[sections={1-3},hideallsubsections]
    \column{0.48\textwidth}
    \tableofcontents[sections={4-6},hideallsubsections]
  \end{columns}
\end{frame}

% ---------------------------------------------------------------------------
% TÍTULO TENTATIVO
% ---------------------------------------------------------------------------
\begin{frame}{Título tentativo de la tesis}
\begin{block}{}
\centering\large
\textbf{Efecto de la heterogeneidad térmica del entorno sobre la temperatura del
conductor en cables eléctricos subterráneos en centrales solares y eólicas: modelado numérico 2D}
\end{block}

\vspace{0.6em}
\begin{columns}[T]
\column{0.32\textwidth}
\textbf{Área disciplinaria}
\begin{itemize}[leftmargin=*]\small
  \item Sistemas eléctricos de potencia
  \item Transferencia de calor computacional
  \item Modelado numérico e IA aplicada
\end{itemize}

\column{0.32\textwidth}
\textbf{Variables centrales}
\begin{itemize}[leftmargin=*]\small
  \item \textbf{X}: conductividad $k(x,y)$ del suelo/relleno
  \item \textbf{Y}: temperatura máxima $T_{\max}$ del conductor
  \item Ampacidad disponible
\end{itemize}

\column{0.32\textwidth}
\textbf{Palabras clave}
\begin{itemize}[leftmargin=*]\small
  \item Underground power cable
  \item Thermal ampacity rating
  \item Soil thermal heterogeneity
  \item Physics-informed neural networks
  \item Dynamic thermal rating (DTR)
\end{itemize}
\end{columns}
\end{frame}

% ---------------------------------------------------------------------------
% CLASIFICACIÓN DE LA INVESTIGACIÓN
% ---------------------------------------------------------------------------
\begin{frame}{Clasificación de la investigación}
\begin{columns}[T]

\column{0.44\textwidth}
\begin{block}{\normalsize Clasificación IEEE}
\scriptsize\textbf{Computing and Processing} /
\textbf{Power, Energy and Industry Applications}
\vspace{0.25em}
\begin{itemize}[leftmargin=*,nosep]\scriptsize
  \item \textbf{Power cables} --- thermal rating, ampacity
  \item \textbf{Heat transfer} --- computational methods
  \item \textbf{Neural networks} --- physics-informed (PINN)
  \item \textbf{Finite element methods} --- 2D heat conduction
  \item \textbf{Soil properties} --- thermal resistivity, drying
\end{itemize}
\end{block}

\vspace{0.4em}
\begin{block}{\normalsize Clasificación ACM CCS}
\scriptsize
\textbf{Computing methodologies} $\to$ Machine learning $\to$ Neural networks\\[2pt]
\textbf{Applied computing} $\to$ Physical sciences and engineering
\end{block}

\column{0.52\textwidth}
\begin{alertblock}{\normalsize Tipo de investigación}
\vspace{0.35em}

\fcolorbox{blue!50!black}{blue!10}{\parbox{0.88\linewidth}{\centering\vspace{3pt}
  {\small\textbf{Investigación Tecnológica}}\\[4pt]
  {\scriptsize Modelado numérico, validación con benchmark,\\
  análisis cuantitativo de la brecha IEC vs.\ FEM}
\vspace{3pt}}}

\vspace{0.45em}

\fcolorbox{green!50!black}{green!10}{\parbox{0.88\linewidth}{\centering\vspace{3pt}
  {\small\textbf{Innovación de Proceso / Producto}}\\[4pt]
  {\scriptsize Herramienta computacional abierta para\\
  cálculo térmico 2D con $k(x,y)$ variable}
\vspace{3pt}}}
\vspace{0.3em}
\end{alertblock}

\vspace{0.4em}
%\begin{exampleblock}{\normalsize Nivel de la investigación}
%\scriptsize
%\textbf{Exploratoria--descriptiva}: caracteriza cuantitativamente la brecha
%entre métodos normativos y el comportamiento bajo condiciones
%heterogéneas de suelo.
%\end{exampleblock}
\end{columns}
\end{frame}


% ===========================================================================
\section{Problemática}
% ===========================================================================

% ---------------------------------------------------------------------------
% INFRAESTRUCTURA CRÍTICA
% ---------------------------------------------------------------------------
\begin{frame}{Cables subterráneos: infraestructura crítica sin margen de error}
\begin{columns}[T]

\column{0.59\textwidth}
\justifying\small
Los cables eléctricos subterráneos son la columna vertebral de la transmisión en
zonas urbanas y la única vía de evacuación de energía en parques solares y eólicos.

\vspace{0.4em}
\textbf{¿Qué ocurre cuando fallan?}
\begin{itemize}[leftmargin=*]\small
  \item Reparación: \textbf{semanas a meses} (excavación, empalmes, pruebas dieléctricas)
  \item En sistemas sin redundancia: corte total de suministro
  \item Impacto en hospitales, agua potable, industria crítica
  \item Para renovables: toda la generación inactiva durante la reparación
  \item Se convierte en \textbf{cuello de botella} para la transición energética
\end{itemize}

\vspace{0.3em}
\refcite{(Fariz et al., 2026, p.~1~[R4]); (Möbius et al., 2025, p.~1~[R5]); (Enescu et al., 2021, p.~2~[R2])}

\column{0.39\textwidth}
\begin{tikzpicture}[>=Stealth,
  box/.style={draw, rounded corners=4pt, align=center, font=\tiny,
              minimum height=0.82cm, text width=2.2cm, inner sep=3pt}]

  \node[box, fill=yellow!30]         (sub)    at ( 0.0, 5.2) {Subestación\\de transmisión};
  \node[box, fill=gray!25]           (cable)  at ( 0.0, 3.8) {Cable\\subterráneo\\(tramo único)};
  \node[box, fill=blue!15]           (ciudad) at (-1.7, 2.3) {Red\\urbana};
  \node[box, fill=green!15]          (solar)  at ( 1.7, 2.3) {Parque\\solar\,/\,eólico};
  \node[box, fill=red!15,
        draw=red!70!black, dashed]   (falla)  at ( 0.0, 0.4)
        {\textbf{FALLA}\\semanas\\sin servicio};

  \draw[->, thick]              (sub.south) -- (cable.north);
  \draw[->, thick] (cable.south west) to[out=240, in=90] (ciudad.north);
  \draw[->, thick] (cable.south east) to[out=300, in=90] (solar.north);
  \draw[->, thick, red!70!black, dashed] (cable.south) -- (falla.north);

  \node[font=\tiny, gray] at (0,-0.3) {Sin redundancia $\Rightarrow$ impacto total};
\end{tikzpicture}
\end{columns}
\end{frame}

% ---------------------------------------------------------------------------
% COSTO DE LA INCERTIDUMBRE
% ---------------------------------------------------------------------------
\begin{frame}{El costo de la incertidumbre térmica}
\begin{columns}[T]
\column{0.48\textwidth}
\begin{alertblock}{Cuando se \textbf{subestima} la temperatura}
\scriptsize
\begin{itemize}[leftmargin=*,itemsep=1pt]
  \item Temperatura real supera el límite del aislamiento XLPE ($T_{\max} = 90\,^{\circ}$C)
  \item Deterioro acelerado; falla prematura no prevista ni diagnosticada
  \item Especialmente grave en cruces urbanos de alta densidad
  \item Consecuencia: reemplazo no planificado, semanas sin servicio
\end{itemize}
\end{alertblock}

\column{0.48\textwidth}
\begin{exampleblock}{Cuando se \textbf{sobreestima} la temperatura}
\scriptsize
\begin{itemize}[leftmargin=*,itemsep=1pt]
  \item Cable opera por debajo de su capacidad real
  \item Sobredimensionamiento innecesario o subutilización del activo
  \item Pérdida de capacidad para integrar más renovables
  \item DTR estático conservador: capacidad dinámica real no aprovechada
  \refcite{(Fariz et al., 2026~[R4])}
\end{itemize}
\end{exampleblock}
\end{columns}

\vspace{0.4em}
\begin{columns}[T]
\column{0.48\textwidth}
\begin{block}{Causa raíz}
\scriptsize Los métodos normativos asumen condiciones \textbf{homogéneas y estacionarias}.
En la realidad, el suelo es heterogéneo, variable con la estación y la historia de carga.
Esta discrepancia es la brecha que se investiga.
\refcite{(IEC 60287-1-1, 2023); (Enescu et al., 2021)}
\end{block}

\column{0.48\textwidth}
\begin{block}{Magnitud de la brecha}
\scriptsize Variaciones de $k_s$ de $0.5$ a $2.5$\,K\,m/W pueden modificar la
ampacidad en \textbf{más del 20\,\%}. El secado térmico puede reducirla
adicionalmente un \textbf{15--30\,\%}.
\refcite{(Khumalo et al., 2025~[R3]); (Anders, 2005~[N3])}
\end{block}
\end{columns}
\end{frame}


% ===========================================================================
\section{Objeto de estudio}
% ===========================================================================

% ---------------------------------------------------------------------------
% ESTRUCTURA INTERNA DEL CABLE
% ---------------------------------------------------------------------------
\begin{frame}{Estructura interna: el cable como sistema de capas concéntricas}
\begin{columns}[T]
\column{0.65\textwidth}
\justifying\small
Un cable de alta tensión está formado por capas concéntricas con funciones
eléctricas y térmicas bien definidas \refcite{(Thue, 2012~[E1]; IEC~60287-1-1, 2023~[N2])}:

\begin{itemize}[leftmargin=*]\scriptsize
  \item \textbf{Conductor} (Cu o Al, 95--1\,200~mm²): fuente primaria de calor
        $Q = I^2 R_\text{ac}$; $k_\text{Cu} \approx 385$~W/(m\,K)
  \item \textbf{Capas semiconductoras}: uniformizan el campo eléctrico;
        $k \approx 0{,}3$~W/(m\,K)
  \item \textbf{Aislamiento XLPE}: barrera térmica principal;
        $k \approx 0{,}286$~W/(m\,K); $T_{\max} = 90\,^{\circ}$C (continuo)
  \item \textbf{Pantalla metálica} (Cu/Al): fuente secundaria por corrientes inducidas;
        factor $\lambda_1$
  \item \textbf{Cubierta exterior} (PE/PVC): protección mecánica;
        $k \approx 0{,}2$--$0{,}4$~W/(m\,K)
\end{itemize}

% \vspace{0.2em}
\begin{alertblock}{}\scriptsize
La resistividad del suelo ($\rho_T = 1/k_s$) es el parámetro de mayor
incertidumbre: variaciones de $0{,}5$ a $2{,}5$\,K\,m/W modifican la
ampacidad en $>$\,20\,\%.
\refcite{(Khumalo et al., 2025~[R3])}
\end{alertblock}

\column{0.35\textwidth}
\centering
\includegraphics[width=0.6\linewidth]{kim2025_cable_xsection}\\[2pt]
{\tiny\color{gray} Sección transversal cable XLPE 154\,kV (7 capas concéntricas).\\
\refcite{Kim et al.\ (2025~[A1]) --- Fig.~2}}
\end{columns}
\end{frame}

% ---------------------------------------------------------------------------
% INSTALACIÓN: ZANJA, BEDDING, BACKFILL
% ---------------------------------------------------------------------------
\begin{frame}{La instalación: zanja, bedding y configuraciones}
\begin{columns}[T]
\column{0.4\textwidth}
\justifying\scriptsize
El cable está integrado en un \textbf{sistema de instalación} cuyo diseño determina
cómo se disipa el calor \refcite{(Ocłoń et al., 2015~[A2]; Kim et al., 2025~[A1])}:

\begin{itemize}[leftmargin=*,nosep]\scriptsize
  \item \textbf{Profundidad} (0{,}8--1{,}5\,m): mayor profundidad
        estabiliza $T_\text{suelo}$ pero alarga la trayectoria de disipación
  \item \textbf{Bedding}: arena fina ($\pm$15\,cm); primera resistencia
        térmica dominante entre cable y suelo nativo
  \item \textbf{Backfill}: relleno de la zanja; tierra excavada (variable)
        o granular seleccionado
  \item \textbf{Thermal backfill (FTB)}: mezcla arena-cemento-agua;
        maximiza $k$ para aumentar ampacidad sin cambiar el cable
  \item \textbf{Configuración}: \textit{trefoil} (menor calor mutuo) vs.\
        \textit{plana} (acceso más fácil)
\end{itemize}

\vspace{0.15em}
{\scriptsize\color{teal!70!black}
\textbf{Nota IEC~60287:} asume $k$ único del medio; no distingue bedding de
suelo nativo, generando error sistemático en instalaciones con materiales
diferenciados. \refcite{(Atoccsa et al., 2024~[A3]; Kim et al., 2025~[A1])}}

\column{0.54\textwidth}
\centering
\includegraphics[width=1.0\linewidth,height=4.0cm,keepaspectratio]{enescu_trench_xsection}\\[1pt]
{\tiny\color{gray} Sección de zanja: cable, bedding, backfill, suelo nativo.\\
\refcite{Enescu et al.\ (2021~[R2]) --- Fig.~1}}

\vspace{0.15em}
\begin{block}{\scriptsize Configuraciones típicas}\scriptsize
\begin{tabular}{ll}
  \textbf{Trefoil} & menor calor mutuo, preferida\\
  \textbf{Plana}   & acceso más fácil, mayor\\
                   & calentamiento mutuo\\
\end{tabular}
\refcite{(Ocłoń et al., 2015~[A2])}
\end{block}
\end{columns}
\end{frame}

% ---------------------------------------------------------------------------
% SUELO VARIABLE + ECUACIÓN DE CALOR
% ---------------------------------------------------------------------------
\begin{frame}{El suelo como medio térmico variable y la física del problema}
\begin{columns}[T]
\column{0.46\textwidth}
\justifying\scriptsize
El suelo \textbf{no es uniforme}: $k_s$ varía con composición mineral, humedad,
compactación y temperatura:

\begin{center}\scriptsize
\renewcommand{\arraystretch}{1.1}
\begin{tabular}{lc}
\toprule
\textbf{Tipo / condición}         & \textbf{$k_s$ (W/(m\,K))} \\
\midrule
Arena seca                         & 0{,}25--0{,}35 \\
Arena húmeda                       & 1{,}0--2{,}0   \\
Arcilla seca                       & 0{,}4--0{,}5   \\
Arcilla saturada                   & 1{,}4--1{,}8   \\
Suelo urbano (relleno mixto)       & 0{,}5--1{,}5   \\
Roca calcárea                      & 2{,}0--3{,}0   \\
Thermal backfill (FTB)             & 1{,}5--3{,}0   \\
\bottomrule
\end{tabular}
\end{center}

\vspace{0.1em}
\begin{alertblock}{\small Secado térmico}\scriptsize
Bajo carga sostenida: $T > 50\,^{\circ}$C cerca del cable $\to$
$k_s$ baja de $\approx$1{,}5 a $< 0{,}5$\,W/(m\,K)
$\to$ $T_\text{cond}$ sube: ciclo de retroalimentación.
Ampacidad se reduce \textbf{20--30\,\%}.
\refcite{(Khumalo et al., 2025, pp.~6--10~[R3])}
\end{alertblock}

\column{0.50\textwidth}
\justifying\scriptsize
\textbf{La ecuación gobernante: difusión de calor 2D}

\begin{block}{\scriptsize Régimen estacionario --- base IEC~60287}
\centering\scriptsize
$\displaystyle
\frac{\partial}{\partial x}\!\Bigl(k\,\frac{\partial T}{\partial x}\Bigr)
+\frac{\partial}{\partial y}\!\Bigl(k\,\frac{\partial T}{\partial y}\Bigr)
+ Q = 0$\\[2pt]
\scriptsize\justifying
Con $k = \mathrm{const}$: solución analítica de Kelvin.\\
Con $k = k(x,y)$: \textbf{no hay solución analítica}
$\Rightarrow$ FEM, FDM o PINNs.
\end{block}

\vspace{0.15em}
\begin{block}{\scriptsize Régimen transitorio --- DTR dinámico}
\centering\scriptsize
$\displaystyle\rho c_p\,\frac{\partial T}{\partial t} =
\nabla\!\cdot\!\bigl(k\,\nabla T\bigr) + Q(t)$\\[2pt]
\scriptsize\justifying
La capacidad térmica $\rho c_p$ controla la respuesta ante ciclos de
carga y sobrecargas (DTR dinámico).
\end{block}

\vspace{0.1em}
\begin{tikzpicture}
  \node[draw=orange!70!black, fill=orange!8, rounded corners=4pt,
        text width=5.8cm, inner sep=4pt, font=\tiny\justifying] {
    \textbf{Contraste de conductividades térmicas:} desde
    $k \approx 0{,}25$\,W/(m\,K) (arena seca) hasta
    $k \approx 385$\,W/(m\,K) (Cu): diferencia de
    \textbf{3 órdenes de magnitud} en el mismo dominio.
    \rcite{(Carslaw \& Jaeger, 1959~[E2]; Aras et al., 2005~[R1])}
  };
\end{tikzpicture}

\end{columns}
\end{frame}

% ---------------------------------------------------------------------------
% TAXONOMÍA DEL OBJETO DE ESTUDIO
% ---------------------------------------------------------------------------
\begin{frame}{Taxonomía del objeto de estudio}
\begin{columns}[T]

\column{0.42\textwidth}
\centering
{\scriptsize\textbf{Estructura jerárquica del objeto de estudio}}\\[0.4em]
\begin{tikzpicture}[>=Stealth, scale=0.88, transform shape,
  root/.style={draw, rounded corners=5pt, fill=teal!55, text=white,
               font=\scriptsize\bfseries, align=center, inner sep=5pt,
               text width=2.3cm, minimum height=0.9cm},
  L1/.style={draw, rounded corners=4pt, font=\tiny\bfseries,
             align=center, inner sep=4pt, text width=2.2cm, minimum height=0.78cm}]

\node[root] (root) at (0, 0) {Cable\\eléctrico\\subterráneo};

\node[L1, fill=orange!30] (tension) at (4.5,  2.4) {Nivel de\\tensión};
\node[L1, fill=blue!22]   (aislam)  at (4.5,  1.2) {Tipo de\\aislamiento};
\node[L1, fill=purple!18] (config)  at (4.5,  0.0) {Configuración\\instalación};
\node[L1, fill=green!22]  (entorno) at (4.5, -1.2) {Entorno\\térmico};
\node[L1, fill=red!18]    (metodo)  at (4.5, -2.4) {Método de\\evaluación};

% trunk line
\draw[gray!50, thick] (1.8, 2.4) -- (1.8, -2.4);
% root → trunk
\draw[->, gray!60, thick] (root.east) -- (1.8, 0);
% trunk → L1
\foreach \y in {2.4, 1.2, 0, -1.2, -2.4} {
  \draw[->, gray!60] (1.8, \y) -- (3.4, \y);
}
% colored dots
\fill[orange!60]  (1.8,  2.4) circle (3pt);
\fill[blue!50]    (1.8,  1.2) circle (3pt);
\fill[purple!50]  (1.8,  0.0) circle (3pt);
\fill[green!50]   (1.8, -1.2) circle (3pt);
\fill[red!50]     (1.8, -2.4) circle (3pt);

\end{tikzpicture}

\column{0.55\textwidth}
\vspace{0.4em}
\setlength{\tabcolsep}{3pt}
\renewcommand{\arraystretch}{0.82}
\begin{tabular}{>{\tiny\bfseries}p{2.0cm}
                >{\tiny}p{2.7cm}
                >{\tiny\color{gray!65}}p{1.3cm}}
\toprule
\textbf{Dimensión} & \textbf{Tipos / variantes} & \textbf{Refs.} \\
\midrule
\rowcolor{orange!18} Nivel de tensión  & MT (1--36\,kV) · AT/EAT & {[N2],[N5]} \\
\rowcolor{blue!12}   Aislamiento       & XLPE · EPR · papel impr. & {[N5],[E1]} \\
\rowcolor{purple!10} Configuración     & Trébol · Plana · ducto & {[N2],[N6]} \\
\rowcolor{green!12}  Entorno térmico   & Suelo nat.\ · backfill · aire & {[A1],[A2]} \\
\rowcolor{red!10}    Método evaluación & IEC\,60287 · FEM · PINN/ML & {[N2],[N6],[R1]} \\
\bottomrule
\end{tabular}

\vspace{0.6em}
\begin{alertblock}{\scriptsize Foco de esta investigación}\scriptsize
Cable \textbf{XLPE MT} en trébol o plano , entorno con \textbf{suelo
heterogéneo} $k(x,y)$.
Evaluación mediante \textbf{FEM 2D} y \textbf{PINN} como alternativa eficiente.
\end{alertblock}
\end{columns}
\end{frame}


% ===========================================================================
\section{Estado del arte}
% ===========================================================================

% ---------------------------------------------------------------------------
% METODOLOGÍA DE BÚSQUEDA --- MINDMAP + STRINGS
% ---------------------------------------------------------------------------
\begin{frame}{Metodología de búsqueda bibliográfica}
\begin{columns}[T]

\column{0.47\textwidth}
\centering
{\scriptsize\textbf{Mapa conceptual de búsqueda}}\\[0.2em]
\begin{tikzpicture}[
  mindmap,
  grow cyclic,
  text=black,
  every node/.style={concept, font=\tiny},
  concept/.append style={fill opacity=0.35, text opacity=1},
  concept color=teal!45,
  level 1/.append style={level distance=2.2cm, sibling angle=90,
                          concept color=blue!40},
  level 2/.append style={level distance=1.4cm, sibling angle=50,
                          concept color=cyan!40, font=\tiny},
  scale=0.67, transform shape]

\node[font=\scriptsize\bfseries, concept color=teal!60]
  {Cables\\eléctricos\\subterráneos\\análisis térmico}
  child [grow=right] { node {Tipo de cable}
    child [grow=north east] { node {\rotatebox{0}{underground cable}} }
    child [grow=south east] { node {\rotatebox{0}{buried cable}} }
  }
  child [grow=up] { node {Aspecto térmico}
    child [grow=north west] { node {\rotatebox{0}{ampacity}} }
    child [grow=north] { node {\rotatebox{0}{thermal rating}} }
    child [grow=north east] { node {\rotatebox{0}{heat transfer}} }
  }
  child [grow=left] { node {Entorno suelo}
    child [grow=south west] { node {\rotatebox{0}{soil resistivity}} }
    child [grow=north west] { node {\rotatebox{0}{backfill}} }
  }
  child [grow=down] { node {Tipo publicación}
    child [grow=south west] { node {\rotatebox{0}{review}} }
    child [grow=south] { node {\rotatebox{0}{assessment}} }
    child [grow=south east] { node {\rotatebox{0}{systematic}} }
  };
\end{tikzpicture}

\vspace{0.3em}
{\tiny\color{gray}
Bases: Scopus, WoS, IEEE Xplore, ScienceDirect.\\
Período: 2000--2026. Idioma: inglés.}

\column{0.5\textwidth}
\begin{block}{\scriptsize String~1 --- Estado del arte}
\tiny\ttfamily
("underground power cable*" OR "buried power cable*"\\
\quad OR "underground cable system*")\\
AND (thermal OR ampacity OR "thermal rating" OR temperature)\\
AND (review OR assessment OR systematic OR survey\\
\quad OR mapping OR overview)
\end{block}

\vspace{0.4em}
\begin{block}{\scriptsize String~2 --- Aportes metodológicos}
\tiny\ttfamily
("underground power cable*" OR "buried cable*")\\
AND (soil OR backfill OR bedding OR "thermal resistivity"\\
\quad OR "soil drying")\\
AND (FEM OR "finite element" OR PINN OR "neural network"\\
\quad OR "machine learning" OR numerical)
\end{block}

\end{columns}
\end{frame}

% ---------------------------------------------------------------------------
% TABLA COMPARATIVA DE REVISIONES
% ---------------------------------------------------------------------------
\begin{frame}{Revisiones sistemáticas y evaluaciones encontradas}
\centering\footnotesize
\setlength{\tabcolsep}{4pt}
\renewcommand{\arraystretch}{1.35}
\begin{tabular}{>{\bfseries\scriptsize}p{2.5cm}
                >{\scriptsize\centering}p{0.6cm}
                >{\scriptsize}p{3.0cm}
                >{\scriptsize}p{2.5cm}
                >{\scriptsize}p{4.0cm}}
\toprule
\textbf{Referencia} & \textbf{Año} & \textbf{Revista / tipo} & \textbf{Cobertura} & \textbf{Hallazgo principal} \\
\midrule
\rowcolor{blue!7}
Aras et al. & 2005 & \textit{EPCS} --- evaluación comparativa & 1957--2005 & IEC~60287 diverge del FEM $>$15\,\% en instalaciones complejas \\
Enescu et al. & 2021 & \textit{Energies} --- revisión sistemática DTR & 2000--2020 & Resistividad del suelo: principal fuente de incertidumbre en DTR \\
\rowcolor{blue!7}
Khumalo et al. & 2025 & \textit{ITEES} --- evaluación crítica & 2000--2024 & Secado del suelo degrada ampacidad 20--30\,\%; IEC~60853-3 insuficiente \\
Fariz et al. & 2026 & \textit{IEEE Access} --- mapeo sistemático & 2015--2025 & 5 brechas abiertas en DTR: suelo heterogéneo, ML inverso, SCADA \\
\rowcolor{blue!7}
Möbius et al. & 2025 & \textit{RSER} --- análisis de impacto & 2010--2024 & Ampacidad varía $\pm$15--20\,\% estacionalmente; ratings estáticos ignoran esto \\
\bottomrule
\end{tabular}

\vspace{0.4em}
{\tiny\color{gray}
EPCS = \textit{Electric Power Comp.\ \& Systems};
ITEES = \textit{Int.\ Trans.\ on Electrical Energy Systems};
RSER = \textit{Renewable and Sustainable Energy Reviews}.}
\end{frame}

% ---------------------------------------------------------------------------
% ENESCU 2021
% ---------------------------------------------------------------------------
\begin{frame}{Enescu et al.\ (2021) --- revisión sistemática DTR}
\begin{columns}[T]
\column{0.65\textwidth}
\begin{block}{\scriptsize Referencia}
\scriptsize Enescu, D., Colella, P., \& Russo, A. (2021). Concepts and methods
to assess the dynamic thermal rating of underground power cables.
\textit{Energies, 14}(9), 2591.
\href{https://doi.org/10.3390/en14092591}{doi:10.3390/en14092591}
\end{block}

\vspace{0.1em}
\begin{block}{\scriptsize Cadena de búsqueda declarada}
\tiny\ttfamily
("underground power cable" OR "buried cable")\\AND ("dynamic thermal rating" OR DTR\\
\quad OR "ampacity")\\AND (method* OR assessment OR review)
\end{block}

\vspace{0.05em}
\textbf{\scriptsize Hallazgos clave}
\begin{itemize}[leftmargin=*,nosep]\tiny
  \item Clasifica métodos DTR: normativos, RC equivalentes, FEM, ML
  \item La \textbf{resistividad del suelo} es la principal fuente de incertidumbre
  \item El monitoreo DTS (fibra óptica) transforma la precisión del DTR
  \item Brecha: falta integración entre modelos físicos y datos operativos
\end{itemize}

\column{0.3\textwidth}
{\scriptsize\textbf{Criterios de inclusión/exclusión}}
\begin{itemize}[leftmargin=*,itemsep=0pt]\tiny
  \item \textbf{Inclusión}: cables subterráneos, métodos de evaluación
        térmica, 2000--2020, Scopus/WoS
  \item \textbf{Exclusión}: cables aéreos, BT, sin validación numérica
\end{itemize}


\begin{alertblock}{\scriptsize Relevancia para esta tesis}\tiny
Confirma que $k(x,y)$ heterogénea es el factor limitante no resuelto en el
DTR. Justifica la pregunta central de la investigación.
\end{alertblock}
\end{columns}
\end{frame}

% ---------------------------------------------------------------------------
% FARIZ 2026
% ---------------------------------------------------------------------------
\begin{frame}{Fariz et al.\ (2026) --- mapeo sistemático DTR}
\begin{columns}[T]
\column{0.46\textwidth}
\begin{block}{\scriptsize Referencia}
\scriptsize Fariz, K., et al.\ (2026). Systematic mapping of dynamic thermal
rating for underground cables. \textit{IEEE Access.}
\href{https://doi.org/10.1109/ACCESS.2025.3650354}{doi:10.1109/ACCESS.2025.3650354}
\end{block}

\vspace{0.15em}
\begin{block}{\scriptsize Cadena de búsqueda declarada}
\tiny\ttfamily
("underground cable" OR "power cable")\\
AND ("thermal rating" OR "ampacity")\\
AND ("dynamic" OR "machine learning"\\
\quad OR "neural network" OR "PINN")
\end{block}

\vspace{0.15em}
\textbf{\small Hallazgos clave}
\begin{itemize}[leftmargin=*,nosep]\scriptsize
  \item Mapeo de 85$+$ estudios (2015--2025) sobre DTR en cables
  \item Clasifica: estadísticos, RC, FEM, PINNs, ML supervisado
  \item \textbf{Problema inverso} ($k$ del suelo desde $T$): sin resolver
  \item Brecha entre modelos físicos y herramientas operativas
\end{itemize}

\column{0.50\textwidth}
\textbf{\small Brechas abiertas (Fariz, 2026)}
\vspace{0.2em}

\begin{tikzpicture}[
  gap/.style={draw, rounded corners=4pt, fill=orange!12,
              text width=6.6cm, font=\scriptsize, inner sep=4pt,
              minimum height=0.62cm},
  >=Stealth]
  \foreach \txt/\y/\col in {
    {B1: Heterogeneidad espacial $k(x,y)$ del suelo}/4.0/red!15,
    {B2: Incertidumbre del backfill y bedding}/3.15/orange!15,
    {B3: Integración con SCADA operacional}/2.3/yellow!20,
    {B4: ML para estimación inversa de $k$}/1.45/green!12,
    {B5: DTR en parques renovables variables}/0.6/blue!10
  }{
    \node[gap, fill=\col] at (2.3,\y) {\txt};
  }
  %\draw[<-, thick, red!60!black] (-0.5,4.0) -- (-0.1,1.45)
   % node[midway, left, font=\tiny\bfseries, rotate=90, color=red!60!black]
    {\hbox{Esta tesis}};
\end{tikzpicture}

\vspace{0.15em}
\begin{alertblock}{\scriptsize Posicionamiento}\scriptsize
Aborda \textbf{B1}: brecha IEC vs.\ modelo con $k(x,y)$ variable.\\
Explora \textbf{B4}: inversión de $k_s$ desde mediciones operativas de $T$.
\end{alertblock}
\end{columns}
\end{frame}

% ---------------------------------------------------------------------------
% KHUMALO 2025
% ---------------------------------------------------------------------------
\begin{frame}{Khumalo et al.\ (2025) --- evaluación crítica: secado del suelo}
\begin{columns}[T]
\column{0.50\textwidth}
\begin{block}{\scriptsize Referencia}
\scriptsize Khumalo, N.~Q., et al. (2025). A critical assessment of cable rating
methods under soil drying out conditions. \textit{Int.\ Trans.\ Electrical Energy
Systems.}
\href{https://doi.org/10.1155/etep/5946564}{doi:10.1155/etep/5946564}
\end{block}

\vspace{0.2em}
\begin{block}{\scriptsize Enfoque de búsqueda}
\tiny\ttfamily
("cable rating" OR "ampacity")\\AND ("soil drying" OR "thermal dryout"\\
\quad OR "moisture migration")\\AND (IEC OR standard OR method)
\end{block}

\vspace{0.15em}
\textbf{\small Hallazgos clave}
\begin{itemize}[leftmargin=*]\scriptsize
  \item El secado del suelo reduce la ampacidad en \textbf{20--30\,\%}
  \item IEC~60853-3 subestima sistemáticamente el riesgo real de secado
  \item Suelos nativos sin backfill más susceptibles al secado
  \item Propone metodología de rating adaptada con medición $k$ in-situ
\end{itemize}

\column{0.46\textwidth}
\centering
\includegraphics[width=0.60\linewidth]{own_kim_fem_zoom}\\[2pt]
{\tiny\color{gray} Cable XLPE 154\,kV: detalle capas concéntricas (zoom FEM).\\
\refcite{Reproducción propia --- Koc y Meléndez, 2025}}

\vspace{0.2em}
\begin{exampleblock}{\scriptsize Implicación clave}\scriptsize
La variabilidad de $k_s$ es tanto \textbf{espacial} como \textbf{temporal}.
Esta dinámica es el corazón del problema investigado.
\end{exampleblock}
\end{columns}
\end{frame}

% ---------------------------------------------------------------------------
% ARAS 2005
% ---------------------------------------------------------------------------
\begin{frame}{Aras et al.\ (2005) --- evaluación comparativa de métodos de ampacidad}
\begin{columns}[T]
\column{0.50\textwidth}
\begin{block}{\scriptsize Referencia}
\scriptsize Aras, F., Oysu, C., \& Yilmaz, G. (2005). An assessment of the
methods for calculating ampacity of underground power cables.
\textit{Electric Power Components and Systems, 33}(12), 1385--1402.
\href{https://doi.org/10.1080/15325000590964425}{doi:10.1080/15325000590964425}
\end{block}

\vspace{0.15em}
\begin{block}{\scriptsize Metodología comparativa}
\scriptsize Compara tres métodos sobre el mismo caso de referencia
(cable 154\,kV en suelo estratificado):\\[2pt]
\begin{tabular}{ll}\scriptsize
  (1) & IEC~60287 analítico \\
  (2) & FEM (elementos finitos) \\
  (3) & BEM (elementos de contorno)\\
\end{tabular}
\end{block}

\vspace{0.15em}


\column{0.46\textwidth}
\centering
\textbf{\small Hallazgos clave}
\begin{itemize}[leftmargin=*]\scriptsize
  \item IEC~60287 diverge del FEM en $\pm$3--8\,\% en condiciones ideales
  \item En suelo estratificado, la divergencia supera \textbf{15\,\%}
  \item FEM captura efectos de $k(x,y)$ variable que IEC ignora
  \item El caso publicado es el \textbf{benchmark} estándar de la literatura
\end{itemize}

%\includegraphics[width=0.90\linewidth]{own_aras_geometry}\\[2pt]
%{\tiny\color{gray} Geometría del caso Aras 2005: 154\,kV trefoil en suelo
%estratificado. Dominio 2D con condiciones de contorno Dirichlet/Neumann.\\
%\refcite{Reproducción propia --- Koc y Meléndez, 2025}}

\vspace{0.3em}
\begin{alertblock}{\scriptsize Benchmark para esta tesis}\scriptsize
El caso Aras 2005 es el caso de validación del modelo computacional propio.
Las diferencias entre IEC y FEM documentadas aquí son la evidencia cuantitativa
de la brecha que se pretende caracterizar y cuantificar sistemáticamente.
\end{alertblock}
\end{columns}
\end{frame}

% ---------------------------------------------------------------------------
% SÍNTESIS ESTADO DEL ARTE
% ---------------------------------------------------------------------------
\begin{frame}{Síntesis del estado del arte: lo que se sabe y lo que falta}
\begin{columns}[T]
\column{0.54\textwidth}
\textbf{\small Lo que el estado del arte \underline{confirma}}\\[0.2em]
\begin{itemize}[leftmargin=*]\small
  \item[$\checkmark$] Métodos normativos divergen del comportamiento real cuando
        el suelo no es uniforme
        \rcite{(Aras, 2005~[R1]; Enescu, 2021~[R2])}
  \item[$\checkmark$] Secado térmico: mecanismo cuantificable que degrada la
        ampacidad 20--30\,\%
        \rcite{(Khumalo, 2025~[R3])}
  \item[$\checkmark$] Variabilidad estacional impacta la ampacidad en $\pm$15--20\,\%
        \rcite{(Möbius, 2025~[R5])}
  \item[$\checkmark$] DTR dinámico mejora el aprovechamiento pero requiere modelos
        más precisos
        \rcite{(Fariz, 2026~[R4])}
  \item[$\checkmark$] FEM/FVM da resultados más precisos pero es complejo y no
        accesible
        \rcite{(CIGRÉ WG B1.56, 2022~[N4])}
\end{itemize}

\column{0.46\textwidth}
\textbf{\small Lo que \underline{todavía falta} (brechas abiertas)}\\[0.2em]
\begin{itemize}[leftmargin=*]\small
  \item[$\times$] Modelo abierto para $k(x,y)$ arbitraria sin
        configuración experta
        \rcite{(Fariz, 2026 --- B1~[R4])}
  \item[$\times$] Herramienta para el \textbf{problema inverso}: estimar
        $k_s$ desde mediciones operativas de $T$
        \rcite{(Fariz, 2026 --- B4~[R4])}
  \item[$\times$] Cuantificación sistemática de la brecha IEC vs.\ real
        en función de parámetros de instalación
  \item[$\times$] Modelos adaptables a suelos nativos (no solo backfill europeo)
        \rcite{(Khumalo, 2025~[R3])}
  \item[$\times$] Integración con sensores SCADA/DTS para DTR en cuasi-tiempo real
\end{itemize}
\end{columns}

%\vspace{0.3em}
\begin{exampleblock}{}
\centering\small
\textbf{Esta investigación se posiciona en las brechas B1 y B4 de Fariz et al.\
(2026)}: caracterizar cuantitativamente la brecha IEC vs.\ modelo detallado
con $k(x,y)$ variable, y explorar la viabilidad del enfoque inverso para
estimar $k_s$ del suelo desde mediciones de temperatura operativa.
\end{exampleblock}
\end{frame}


% ===========================================================================
\section{Formulación del problema}
% ===========================================================================

% ---------------------------------------------------------------------------
% GAP DIAGRAM
% ---------------------------------------------------------------------------
\begin{frame}{La brecha: de dónde venimos y a dónde queremos llegar}
\begin{center}
\resizebox{0.96\textwidth}{!}{%
\begin{tikzpicture}[>=Stealth,
  estado/.style={draw, rounded corners=5pt, text width=3.8cm, align=center,
                 minimum height=3.2cm, font=\small, inner sep=7pt},
  lbl/.style={font=\tiny\bfseries, text width=3.8cm, align=center}]

  \node[estado, fill=red!10, draw=red!55!black] (actual) at (0,0) {
    \textbf{Estado actual}\\[5pt]
    \footnotesize
    Normas IEC / Neher-McGrath\\con supuestos uniformes\\[3pt]
    FEM preciso pero no\\accesible ni adaptable\\[3pt]
    Sin inversión de\\parámetros del suelo\\[3pt]
    Incertidumbre en $T_\text{cond}$\\no cuantificada
  };

  \node[estado, fill=orange!14, draw=orange!60!black] (brecha) at (5.6,0) {
    \textbf{Brecha}\\[5pt]
    \footnotesize
    $k(x,y)$ heterogénea\\ignorada o asumida cte.\\[3pt]
    Casos complejos fuera\\de premisas normativas\\[3pt]
    Secado, ciclos y\\clima no integrados\\[3pt]
    Sin herramienta abierta\\y adaptable
  };

  \node[estado, fill=green!8, draw=green!50!black] (deseable) at (11.2,0) {
    \textbf{Estado deseable}\\[5pt]
    \footnotesize
    Modelo 2D con $k(x,y)$\\variable por zonas\\[3pt]
    Problema directo\\e inverso resueltos\\[3pt]
    Open-source, adaptable\\y validado\\[3pt]
    Eficiente para DTR\\y mantenimiento
  };

  \draw[->, very thick, red!70!black]
    (actual.east) -- node[above, font=\tiny, color=red!70!black]
    {divergencia} (brecha.west);

  \draw[->, very thick, dashed, green!55!black]
    (brecha.east) -- node[above, font=\tiny, color=green!50!black]
    {cerrar} (deseable.west);

  \node[lbl, below=0.35cm of actual, text=red!70!black]
    {Métodos insuficientes\\para condiciones reales};
  \node[lbl, below=0.35cm of brecha, text=orange!80!black]
    {Problema a resolver\\(esta investigación)};
  \node[lbl, below=0.35cm of deseable, text=green!50!black]
    {Meta: precisión,\\seguridad y eficiencia};
\end{tikzpicture}%
}
\end{center}
\end{frame}

% ---------------------------------------------------------------------------
% FORMULACIÓN DEL PROBLEMA
% ---------------------------------------------------------------------------
\begin{frame}{Formulación del problema de investigación}
\begin{exampleblock}{Síntesis}
\justifying\small
Los métodos predominantes de evaluación térmica de cables eléctricos
subterráneos ---basados en normas estacionarias (IEC~60287, IEC~60853,
Neher-McGrath) o en modelos FEM de uso restringido--- asumen condiciones
térmicas homogéneas y no están diseñados para representar la
\textbf{heterogeneidad real del suelo}, los casos de instalación complejos
ni la estimación inversa de propiedades térmicas desde mediciones operativas.

\vspace{0.3em}
Esta limitación genera una \textbf{brecha} entre la temperatura estimada del
conductor y la temperatura real bajo condiciones de suelo variable, secado
térmico, cargas dinámicas o arreglos de múltiples circuitos.

\vspace{0.3em}
La brecha se expresa como: (1)~incertidumbre sobre la ampacidad disponible,
(2)~riesgo de sobrecalentamiento no detectado, y (3)~falta de herramientas
accesibles para el diseño y operación de instalaciones fuera del caso base
normativo. Cerrar esta brecha permitiría diseñar instalaciones más seguras,
operar con mayor aprovechamiento de la capacidad térmica real y habilitar el
mantenimiento predictivo de activos críticos sin redundancia.
\end{exampleblock}

\vspace{0.3em}
\refcite{(Neher y McGrath, 1957~[N1]); (IEC~60287-1-1, 2023~[N2]); (Aras et al., 2005~[R1]);
(Khumalo et al., 2025~[R3]); (Enescu et al., 2021~[R2]); (Fariz et al., 2026~[R4])}
\end{frame}


% ===========================================================================
\section{Preguntas y objetivos}
% ===========================================================================

% ---------------------------------------------------------------------------
% PREGUNTA CENTRAL + SUBORDINADAS
% ---------------------------------------------------------------------------
\begin{frame}{Pregunta de investigación}

\begin{block}{Pregunta central}
\justifying
¿En qué medida los métodos simplificados de evaluación térmica de cables
eléctricos subterráneos representan adecuadamente la distribución de temperatura
real cuando el suelo presenta conductividad térmica espacialmente variable
$k(x,y)$, y cuál es el efecto cuantificable de dicha variabilidad sobre la
temperatura máxima del conductor y la ampacidad estimada?
\end{block}

\vspace{0.6em}
\textbf{Preguntas subordinadas:}
\begin{itemize}[leftmargin=*]\small
  \item ¿Cuáles son los supuestos de IEC~60287 y del método Neher-McGrath que
        limitan su validez bajo condiciones reales de suelo e instalación?
  \item ¿Cuánto varía la temperatura del conductor al cambiar la distribución
        espacial $k(x,y)$ del suelo y del material de relleno?
  \item ¿Qué configuraciones de instalación se alejan más de las premisas normativas
        y generan mayor incertidumbre térmica?
  \item ¿Qué información operativa y qué tipo de modelo serían necesarios para
        estimar la ampacidad dinámica real con mayor precisión?
\end{itemize}
\end{frame}

% ---------------------------------------------------------------------------
% OBJETIVOS
% ---------------------------------------------------------------------------
\begin{frame}{Objetivos de la investigación}

\begin{block}{Objetivo general}
\justifying\small
Caracterizar y cuantificar la brecha entre los métodos normativos simplificados y
el comportamiento térmico real de cables eléctricos subterráneos, identificando
las condiciones de suelo, geometría de instalación y perfil de carga que
determinan la magnitud de dicha brecha y los requisitos para cerrarla.
\end{block}

\vspace{0.4em}
\textbf{Objetivos específicos:}
\begin{enumerate}[leftmargin=*]\small
  \item Describir y sistematizar los supuestos y limitaciones de los métodos
        normativos vigentes (IEC~60287, IEC~60853, Neher-McGrath) para el
        cálculo de ampacidad en condiciones no ideales.
  \item Identificar los casos de instalación o condición de entorno donde los
        métodos simplificados presentan mayor desviación respecto a modelos
        numéricos más detallados.
  \item Analizar cuantitativamente el efecto de la variabilidad de $k(x,y)$
        del suelo/backfill sobre $T_{\max}$ del conductor y la ampacidad estimada.
  \item Delimitar los requisitos funcionales de un modelo computacional abierto
        para representar adecuadamente la transferencia de calor 2D con $k(x,y)$
        variable y condiciones de instalación complejas.
\end{enumerate}
\end{frame}


% ===========================================================================
\section{Referencias}
% ===========================================================================

\begin{frame}{Referencias}
\begin{columns}[T]
\column{0.48\textwidth}

{\footnotesize\color{blue!60!black}\textbf{[N] Normas y libros de referencia}}
\vspace{0.1pt}
\begin{itemize}[leftmargin=*,itemsep=1pt]\scriptsize
  \item \textbf{[N1]} Neher \& McGrath (1957). Calc.\ of temperature rise and
        load capability of cable systems. \textit{AIEE Trans.}
  \item \textbf{[N2]} IEC~60287-1-1:2023. \textit{Calculation of current rating.}
  \item \textbf{[N3]} Anders (2005). \textit{Rating of power cables in unfavorable
        thermal environment.} IEEE Press.
  \item \textbf{[N4]} CIGRÉ WG B1.56 (2022). \textit{Cable rating examples for
        calculation tool verification.}
  \item \textbf{[N5]} IEC~60502-2:2014. \textit{Power cables with extruded
        insulation, 1\,kV -- 30\,kV.} IEC.
  \item \textbf{[N6]} CIGRÉ WG B1.87 (2025). \textit{Finite element analysis
        for cable rating calculations.} TB~963.
\end{itemize}

\vspace{0.1pt}
{\footnotesize\color{teal!70!black}\textbf{[R] Revisiones y evaluaciones (estado del arte)}}
\vspace{0.1pt}
\begin{itemize}[leftmargin=*,itemsep=1pt]\scriptsize
  \item \textbf{[R1]} Aras et al.\ (2005). Assessment of methods for calculating
        ampacity. \textit{EPCS, 33}(12). \href{https://doi.org/10.1080/15325000590964425}{doi}
  \item \textbf{[R2]} Enescu et al.\ (2021). Concepts and methods for DTR of
        underground cables. \textit{Energies, 14}(9). \href{https://doi.org/10.3390/en14092591}{doi}
  \item \textbf{[R3]} Khumalo et al.\ (2025). Critical assessment of cable
        rating under soil drying. \textit{ITEES.} \href{https://doi.org/10.1155/etep/5946564}{doi}
  \item \textbf{[R4]} Fariz et al.\ (2026). Systematic mapping of DTR for
        underground cables. \textit{IEEE Access.} \href{https://doi.org/10.1109/ACCESS.2025.3650354}{doi}
  \item \textbf{[R5]} Möbius et al.\ (2025). Cable ampacity: seasonal and
        climate-related changes. \textit{RSER.}
\end{itemize}

\column{0.48\textwidth}

{\footnotesize\color{orange!70!black}\textbf{[A] Aportes metodológicos}}
\vspace{0.1pt}
\begin{itemize}[leftmargin=*,itemsep=1pt]\scriptsize
  \item \textbf{[A1]} Kim et al.\ (2025). Heat transfer in underground cable
        with new bedding material. \textit{Geothermics, 125}.
  \item \textbf{[A2]} Ocłoń et al.\ (2015). Optimizing underground cable
        bedding using PSO. \textit{Energy, 92}(P2).
  \item \textbf{[A3]} Atoccsa et al.\ (2024). Ampacity optimization with
        thermal backfill using dynamic PSO. \textit{Energies, 17}(5).
\end{itemize}

\vspace{0.1pt}
{\footnotesize\color{gray!70!black}\textbf{[E] Textos clásicos de soporte}}
\vspace{0.1pt}
\begin{itemize}[leftmargin=*,itemsep=1pt]\scriptsize
  \item \textbf{[E1]} Thue, W.~A. (2012). \textit{Electrical Power Cable
        Engineering.} CRC Press.
  \item \textbf{[E2]} Carslaw \& Jaeger (1959). \textit{Conduction of Heat
        in Solids.} Oxford UP.
\end{itemize}

\vspace{0.1pt}
\begin{block}{}\scriptsize
\textbf{Código}: \textcolor{blue!60!black}{[N]} norma/libro,
\textcolor{teal!70!black}{[R]} revisión/evaluación (estado del arte),
\textcolor{orange!70!black}{[A]} aporte metodológico,
\textcolor{gray!70!black}{[E]} texto clásico de soporte.
\end{block}

\end{columns}
\end{frame}

% ---------------------------------------------------------------------------
% CIERRE
% ---------------------------------------------------------------------------
\begin{frame}
\centering
{\Large \textbf{Gracias}}

\vspace{0.8em}
\large Luis Koc \quad $|$ \quad Herbert Meléndez

\vspace{0.4em}
\normalsize Trabajo de Investigación I $\cdot$ MIA 3er Ciclo A 2026-1 $\cdot$ FIIS/UNMSM

\vspace{0.6em}
\begin{block}{\centering Síntesis del problema}
\centering\small
Existe una brecha entre los métodos simplificados de evaluación térmica de cables
eléctricos subterráneos y la realidad del sistema bajo condiciones heterogéneas,
dinámicas y complejas. El estado del arte (Enescu 2021, Khumalo 2025, Fariz 2026)
confirma la brecha pero carece de herramientas abiertas para cuantificarla
sistemáticamente. Esta investigación apunta a cerrar esa brecha mediante modelos
numéricos 2D con $k(x,y)$ variable y enfoque directo e inverso.
\end{block}
\end{frame}

\end{document}
